\section{Conclusions and Future Work}
\label{chap:conclusions}

In this work, we explored the use of autoencoders for speckle noise reduction in SAR images, presenting original contributions in both the training data generation methodology and the development of flexible tools for experimenting with different architectures. We demonstrated that it is possible to use the $\mathcal{G}_I^0$ distribution to generate synthetic training data for training autoencoders capable of filtering speckle noise, although with limitations in their generalization to actual images. Our experiments consistently reveal that deeper and more complex architectures yield better results. And we also achieved results similar to those of traditional filters. However, the superior performance observed with models trained directly on actual data underscores the importance of having training data that is representative of the application domain.

Our findings highlight a key dichotomy in the use of synthetic data: the $\mathcal{G}_I^0$ model offers a controllable environment for architecture development, but a gap remains with real SAR images. Thus, the strongest use of synthetic data is as a pre-training step to provide a statistical prior, followed by fine-tuning on a smaller real SAR dataset. This two-stage strategy can combine synthetic fidelity with real data adaptation and should be validated using both reconstruction error and speckle-preserving metrics.

Future research will focus on exploring more sophisticated architectures, such as U-Net and attention-based models, and on formalizing the hybrid training strategy suggested by our results. A promising next step is to evaluate transfer learning by pre-training on large synthetic $\mathcal{G}_I^0$ datasets and then fine-tuning on smaller real SAR sets. This may involve layered fine-tuning, mixed-domain batches, and hybrid losses that balance reconstruction accuracy with statistical consistency. We also plan to refine synthetic generation with greater spatial diversity and richer texture variations. In addition, although this work compares the proposed autoencoder approach with the classical Lee filter, a more complete evaluation should include modern speckle reduction algorithms such as the FANS filter. A rigorous comparison with other state-of-the-art methods remains a priority for future work. Finally, we will extend the proposed methodology to more complex data, including polarimetric and multi-temporal SAR images, to assess its performance in a broader range of applications. While the use of Mean Squared Error (MSE) as the loss function provided a solid baseline for this study, we acknowledge its limitations in capturing the perceptual and textural nuances of SAR imagery. As noted, optimizing for MSE can lead to a 'regression to the mean,' resulting in blurred reconstructions that may lack high-frequency details. Our decision was pragmatic, aimed at first validating our data generation and architectural hypotheses in a controlled manner. Thus, our work separates the optimization objective (MSE) from the evaluation criteria: MSE is used to train stable reconstructions, while the speckle-specific metrics verify that the output maintains the desired statistical behavior.

For future work, we plan to explore more advanced loss functions. A promising direction is the development of a custom perceptual loss tailored for SAR images, which would require pre-training a feature extractor on a large SAR dataset. This type of loss compares learned features rather than individual pixels, helping to preserve textural and structural details that MSE alone tends to smooth.

Another direction is to incorporate terms that measure the match between distributions, such as those used in Generative Adversarial Networks (GANs) or Wasserstein-based losses. Implementing this would require a discriminative component or GAN-style architecture capable of learning the statistical distribution of real clean SAR images, which is more complex than our current MSE-based framework.

The positive results from our current MSE-based models provide a strong foundation for these future explorations, which we believe will yield significant improvements in reconstruction quality.
In conclusion, this work establishes a solid foundation for using autoencoders for speckle noise reduction in SAR images, demonstrating the potential of combining statistical models with deep learning to address classic radar image processing problems.